\documentclass[12pt]{article}
\usepackage{color}
\usepackage{amsmath}
\usepackage{arxiv}
\usepackage{lyu}
\usepackage{xcolor}
\usepackage{pgfplots}
\usetikzlibrary{calc}
\pgfplotsset{compat=1.18}
\pgfplotsset{colormap={violet}{rgb255=(25,25,122) rgb255=(238,140,238) color=(white)}}
\usepackage[utf8]{inputenc} % allow utf-8 input
\usepackage[T1]{fontenc}    % use 8-bit T1 fonts
\usepackage{hyperref}       % hyperlinks
\usepackage{url}            % simple URL typesetting
\usepackage{booktabs}       % professional-quality tables
\mathchardef\mhyphen="2D
\usepackage{amsfonts}       % blackboard math symbols
\usepackage{nicefrac}       % compact symbols for 1/2, etc.

\usepackage{microtype}      % microtypography
\usepackage{lipsum}
\usepackage{comment}
\usepackage{subcaption}
\usepackage{graphicx}
\usepackage{soul}           % highlight text \hl{...} and strike \st{...}
\definecolor{myGreen}{rgb}{0.17254902, 0.62745098, 0.17254902}
\newcommand{\francesca}[1]{\color{red}#1 -FG \color{black}}
\newcommand{\huan}[1]{\color{myGreen}\ #1 -HL \color{black}}
\newcommand{\lyu}[1]{\color{red}\ #1 -LL \color{black}}
\DeclareMathOperator{\sign}{sign}
\usepackage{algorithm}
\usepackage{algpseudocode}
\usepackage{appendix}
\usepackage{multirow}
\usepackage{amssymb}
\usepackage{amsthm}
\usepackage{amssymb}
\usepackage{mathabx}
\usepackage{tikz}
\usetikzlibrary{shapes,arrows}
\usetikzlibrary{positioning}
\usetikzlibrary{arrows,
    chains,
    decorations.markings,
    shadows, shapes.arrows,shapes, fit}
\newtheorem{theorem}{Theorem}[section]
\newtheorem{lemma}[theorem]{Lemma}
\newtheorem{proposition}[theorem]{Proposition}
\newtheorem{corollary}[theorem]{Corollary}
\newtheorem{definition}[theorem]{Definition}
\newtheorem{remark}[theorem]{Remark}
\title{Stochastic Branching method for solving the Keller–Segel chemotaxis Model}
\renewcommand{\shorttitle}{}
\author{
  Liyao Lyu\\
  Department of Computational Mathematics, Science \& Engineering, Michigan State University, MI 48824, USA
   \And
  Huan Lei\\
  Department of Computational Mathematics, Science \& Engineering, Michigan State University, MI 48824, USA\\
  Department of Statistics \& Probability, Michigan State University, MI 48824, USA}
\begin{document}
\maketitle

% keywords can be removed
\keywords{} 
\begin{abstract}
We propose an efficient stochastic branching method for solving parabolic-parabolic Keller-Segel (KS) chemotaxis systems. In this approach, both the cell density and the chemical concentration are modeled as empirical measures of particles. By approximating the original problem linearly through a Kolmogorov forward equation for the particles at each time step, our method provides an effective solution to a system of advection-diffusion-reaction equations. Compared to traditional methods, our approach offers better handling of singularities. We present numerical examples in two-dimensional space to demonstrate the accuracy and effectiveness of the proposed method.

\end{abstract}


\section{Introduction}
Chemotaxis models are widely used to describe phenomena such as immune cell migration, tumor growth, and developmental processes. Understanding and solving these models can lead to insights into disease progression and potential therapeutic strategies. A particular challenge arises in resolving the singular behavior associated with the concentration of $u$ at a Dirac mass. Traditional methods, often based on mesh-based discretization, are unable to efficiently represent such singularities without introducing significant errors or instabilities in the solution.

In this paper, we study the Keller-Segel (KS) chemotaxis model \cite{keller1970initiation, patlak1953random} in two spatial dimensions, focusing on the following widely studied formulation:
\begin{equation}\label{equ:KS}
\begin{aligned}
\partial_t u - \nabla \cdot ( \nabla u - \chi(u) \nabla v) &= f_u(u,v), & \bx\in \mathbb{R}^2 , t\in(0,T), \\
v_t - \Delta v &= f_v(u,v), & \bx\in \mathbb{R}^2 , t\in(0,T),
\end{aligned}
\end{equation}

where $u$ and $v$ represent the cell density and chemical concentration, respectively. Chemotaxis is a process in which cells navigate toward areas with favorable chemical concentrations and away from unfavorable ones. The nonlinearity of this process leads to complex behavior, such as blow-up phenomena, where the cell density $u$ can evolve into a Dirac delta function in finite time under certain initial conditions \cite{filbet2006finite, gajewski1998global, herrero1997finite, nagai1995blow}. This singular behavior presents significant challenges for traditional numerical methods, such as finite element \cite{nakaguchi2002fully, saito2011error, strehl2013positivity}, finite volume \cite{filbet2006finite}, and discontinuous Galerkin methods \cite{epshteyn2009new}, which are difficult to accurately capture such solutions in regular functional spaces.



To address this, one natural idea is to approximate the Dirac mass using particle-based methods. However, existing particle methods, such as the partial in-cell method \cite{birdsall2018plasma, hockney2021computer} for solving the Vlasov-Poisson equation, are often based on the assumption that the underlying physical laws are well understood. In biological systems like chemotaxis, where these governing laws are poorly known or difficult to model, such methods are less applicable. Other methods, such as stochastic particle approximation \cite{havskovec2009stochastic, wang2025novel}, rely on interacting particle systems that converge to the original PDE. However, there remains a lack of a unified framework for discretizing nonlinear, nonconservative PDEs into particle-based methods.
While recent advances using the Feynman-Kac formula \cite{han2024deep, hutzenthaler2021speed, nguwi2024deep} and backward stochastic differential equations \cite{han2017deep, han2018solving, hure2020deep} have shown promise for solving a broad class of PDEs, they often introduce an optimization phase that adds complexity and can hinder both the accuracy and efficiency of the method.

In this work, we propose a \textit{stochastic branching method} for solving the Keller-Segel chemotaxis model, which addresses these challenges. Our method approximates the original problem through a linearized Kolmogorov forward equation at each time step, with the solution represented by the density of a branching stochastic process. The stochastic branching method is well-suited for chemotaxis due to its flexibility in representing particle interactions and singularities inherent in the system. This method allows for efficient tracking of individual particles, which is crucial for accurately capturing the blow-up behavior that often occurs in chemotaxis systems.
Although we focus on the 2D Keller-Segel model, the proposed framework is general and can be applied to other reaction-diffusion systems, such as the Gray-Scott model, Allen-Cahn equation, and Schnakenberg pattern formation. Moreover, by leveraging recent advances in density estimation techniques \cite{novikov2021tensor, tang2025wavelet, sheng2025approximation, tang2024solving}, our method can be readily extended to higher-dimensional systems.

In summary, we propose a novel stochastic branching method for solving the Keller-Segel chemotaxis model, overcoming the limitations of traditional numerical methods by efficiently handling singularities and nonlinearities. We demonstrate the effectiveness of our approach through several numerical examples and discuss its extension to higher-dimensional systems.



\section{Stochastic branching method}
The core idea behind our stochastic branching method for solving PDEs is to approximate the solution using the probability density $\rho (t, \bx)$ associated with an empirical measure $\mu(t, \intd \bx) = \rho (t, \bx)\intd \bx$. The empirical measure is then approximated by a set of particles $\bX^i(t)$, for $t\in[0,T]$ and $i=1\cdots N(t)$. By increasing the number of particles $N(0)\to \infty$, we aim for the empirical measure to converge to the true solution   
\[\lim_{N\to \infty} \mu^N_t(\intd \bx) = \lim_{N(0)\to \infty} \frac{1}{N(0)}\sum_{i=1}^{N(t)} \delta_{\bX^i(t)}=\mu(t,\intd \bx).
\]  
It is important to note that the equation we consider here is non-conservative. As a result, the total mass of the empirical measure $\mu(t,\intd \bx)$ may vary over time, and $\rho(t,x)$ is not normalized for all $t$. To explain our main idea, i we simplify the original PDE system (Equation \eqref{equ:KS}) by first considering a scalar-valued reaction-diffusion equation of the form
\begin{equation}\label{equ:reaction-diffusion}
\partial_t u = \nabla \cdot \left( a(u)u + \nabla (d(u) u) \right) + r(u)u.
\end{equation}
Although we initially focus on this scalar case, we note that generalizing the method to more complex systems, such as the one in Equation~\eqref{equ:KS}, follows in a straightforward manner.
\subsection{Initial condition and Normalization}
To connect the solution $u(t,\bx)$ to the probability density $\rho(t,\bx)$, we define a linear transform between them. Since the inital density $\rho(0,\bx)$ is normalized always, we define 
\begin{equation}
    u(t,\bx) = Z \rho(t,\bx),
\end{equation}
where $Z = \int_{\Omega} u(0,\bx)\intd \bx$. 
Consequently, the initial particle distribution should follow the normalized density $\rho(0,\bx)= \frac{u(0,\bx)}{Z}$. At initial time $t=0$, we sample $N_0=N(0)$ particles $\bX_0^i$ from this distribution using the Metropolis–Hastings algorithm. As a result, the equation of $\rho$ satisfies 
\begin{equation}\label{equ:normalized_diffusion_reaction}
\partial_t \rho = \nabla \cdot \left( \tilde a(\rho)\rho + \nabla (\tilde 
 d(\rho) \rho )\right) + \tilde r(\rho)\rho,
\end{equation}
 where $\tilde a(\rho)  = a(Z\rho)$, $\tilde d(\rho)  = d(Z\rho)$, and $\tilde r(\rho)  = r(Z\rho)$. 

\subsection{Operator Splitting}
Next, we introduce an operator splitting strategy that separates the evolution dynamics based on their underlying physical mechanisms. 
% For simplicity, we first consider a scalar-valued reaction-diffusion equation of the form
% \begin{equation}\label{equ:reaction-diffusion}
% \partial_t u = \nabla \cdot \left( a(u)u + d(u)\nabla u \right) + r(u)u,
% \end{equation}
% where the first term models advection and diffusion, and the function $c(u)$ accounts for local reaction dynamics. 
Equation~\eqref{equ:normalized_diffusion_reaction} can be decomposed into two components:
\begin{equation}
\partial_t \rho  = \mathcal A \rho + \mathcal B \rho ,
\end{equation}
where $\mathcal A$ captures the advection-diffusion contributions defined as 
\begin{equation}
 \mathcal A \rho =  \nabla \cdot \left( \tilde a(\rho)\rho + \nabla (\tilde d(\rho) \rho) \right),
\end{equation}
and $\mathcal B\rho$ represents the reaction dynamics, such as cell proliferation and death, given by
\begin{equation}
\mathcal B \rho = \tilde r(\rho)\rho,
\end{equation}
represent the effects of reaction dynamics or cell proliferation and death.
Based on the splitting, we can approximate the solution by 
\begin{equation}
    \rho(t_{n+1})=  \exp\left((\mathcal B+\mathcal A)\tau \right)\rho(t_{n}),
\end{equation}
where $\tau = t_{n+1}-t_n$.
To isolate the contributions of each physical process, we apply a first-order Lie-Trotter splitting approximation:
\begin{equation}
\rho_{n+1} :=\exp(\tau \mathcal B)\exp(\tau \mathcal A)\rho_n,
\end{equation}
where $\rho_n$ is the numerical approximation of $\rho(t_{n})$. 
This first-order splitting provides a simple and efficient approximation, though higher-order splitting methods may be considered in future work for improved accuracy. This decomposition allows us to design particle-based schemes that target each sub-process individually, i.e.
\begin{equation}
\rho_{n+1} =\exp(\tau \mathcal B)\rho^*_n, \quad \rho^*_n = \exp(\tau \mathcal A)\rho_n.
\end{equation}

\subsection{Advection Diffusion Step}
In this step, we model the movement of particles, where each particle follows a stochastic path determined by both a drift term and a diffusion term. At time step $n$, assume we have a set of particles 
$\bX^i_n$, for $i=1,\cdots N_n$, such that the empirical measure 
\[
\mu_n= \frac{1}{N_0}\sum_{i=1}^{N_n}\delta_{\bX_n^i}
\] 
approximates the density, i.e.,  $\mu_n(\intd \bx) = \rho_n(\bx)\intd \bx$. 
Our goal is to evolve these particles according to the advection-diffusion transition probability $p(\bx,s;\by,t)$, which leads to an updated empirical measure
\[\mu_n^* = \int p(\bx,t_n;\by,t_n+\tau)  \mu_n(\intd \bx) =   \frac{1}{N_0}\sum_{i=1}^{N_n}\delta_{\bX_{n,*}^i}
\]  
reflecting the action of the advection-diffusion operator on the density.
By the Kolmogorov forward equation, the transition density $p(\bx,s;\by,t)$ satisfies 
\begin{equation}
    \partial_t p = \tilde {\mathcal A }\rho := - \nabla \cdot(b  p) + \frac{1}{2}\nabla^2:(\sigma^T\sigma p) 
\end{equation}
with initial condition $\lim_{t\to s^+} p(\bx,s;\cdot,t) = \delta_\bx$. By matching the differential operator $\mathcal A$ in our PDE, we have 
\begin{equation}
    b = -\tilde a, \quad \frac{1}{2}\sigma^T \sigma = \tilde d.
\end{equation}
Therefore, the evolution of each particle $\bX^i_n$ as a stochastic process satisfying the It\^o stochastic differential equation
\begin{equation}\label{equ:overd_langevin}
    \intd \bX_n^i = - b(\bX_n^i) \intd t  + \sigma(\bX_n^i) \intd \bW_t,
\end{equation}
where $\bW_t$ is a $d\mhyphen$dimensional Brownian motion, $b$ and $\sigma$ are Lipshitz continuous functions.
% This implies that  the time  evolution of the empirical measure under the SDE is governed by 
% \begin{equation}
% \partial_t \int f(\mathbf{y}) \mu(t, \mathrm{d}\mathbf{y}) = \int f(\mathbf{y}) \tilde {\mathcal{A}}^\dagger \mu(t, \mathrm{d}\mathbf{y}),
% \end{equation}
% where the adjoint Fokker–Planck operator is $\tilde {\mathcal A}  (\rho) =   - \nabla \cdot(b  p) + \frac{1}{2}\nabla^2:(\sigma^T\sigma p) $.  Hence, the empirical measure after the advection-diffusion step satisfies
% \begin{equation}
%     \mu^*_n = \frac{1}{N_0}\sum_{i=1}^{N_n} \delta_{\mathbf{X}^i_{n,*}}  = \frac{1}{N_0} \sum_{i=1}^{N_n} \exp(\tilde {\mathcal A}\tau)\delta_{\mathbf{X}^i_{n}}   =  \exp(\tilde {\mathcal A}\tau) \frac{1}{N_0} \sum_{i=1}^{N_n} \delta_{\mathbf{X}^i_{n}} = \exp(\tilde {\mathcal A}\tau)  \mu_n.
% \end{equation}


\subsection{Brancing Step}
This step modifies the mass of the empirical measure and may change the number of particles over time. Given the intermediate empirical measure from the advection-diffusion step $\mu_n^*$, we aim to evolve this measure under the action of the reaction operator $\mathcal  B$ over a time step $\tau$, such that the resulting empirical measure
\begin{equation}
    \mu_{n+1} = \frac{1}{N_0} \sum_{i=1}^{N_{n+1}} \delta \bX^i_{n+1}
\end{equation}
approximates the updated density $\mu_{n+1}(\intd \bx ) = \rho_{n+1}(\bx) \intd \bx = \exp(\mathcal B\tau) \rho_{n}^*(\bx) \intd \bx = \exp(\mathcal  B\tau) \mu_{n}^*(\intd \bx)$. Therefore, the particles in a small neighbour of $\bx$, i.e., $\bX^i_{n,*}\in \Omega_\epsilon(\bx) := \{\by:|\by-\bx|<\epsilon\}$ , should birth and death with probability $\tilde r(\bx)$ depending on the positiveness of $\tilde r$.
To implement this step, we adopt a stochastic branching process. For each particle $\bX^i_{n,*}$, we associate a branching (or survival) rate determined by the function $\tilde r $. Specifically, over the interval $[t_n,t_{n+1}]$, each particle undergoes one of the following stochastic events:
\begin{itemize}
    \item Survival (no branching): particle $\bX^i_{n,*}$ survival with probability approximately $\exp(-|\tilde r(\bX^i_{n,*})|\tau)$
    \item Death (removal): if $\tilde r<0 $, particle $\bX^i_{n,*}$ dead  with probability $1-\exp(\tilde r(\bX^i_{n,*})\tau)$
    \item Duplication (birth): if $\tilde r>0 $, particle $\bX^i_{n,*}$ duplication  with probability $1-\exp(-\tilde r(\bX^i_{n,*})\tau)$
\end{itemize}

This process realizes the operator $\exp(\mathcal B\tau)$ at the particle level. The expected mass evolution of the empirical measure is consistent with the solution of the ODE:
\begin{equation}
\partial_t \rho = \tilde{r}(\rho) \rho,
\end{equation}
which governs exponential-type growth or decay, depending on the sign and form of $\tilde r$



\subsection{Nonlinear Interpolation}
To evaluate the functions $\tilde a(\rho)$ ,  $\tilde d(\rho)$ , and $\tilde r(\rho)$ , that govern the drift, diffusion, and reaction dynamics in our particle-based scheme, we must first reconstruct the continuous density 
$\rho$ from its empirical approximation. Specifically, given an empirical measure
\[
\mu_n= \frac{1}{N_0}\sum_{i=1}^{N_n}\delta_{\bX_n^i},
\] 
we seek a smooth approximation of $\rho_n$ that allows pointwise evaluation and differentiation as required by the operators.

In this work, we choose to interpolate the empirical distribution into the Fourier space. That is, we project the density onto a finite set of Fourier basis functions, allowing for efficient representation and differentiation. The Fourier coefficients are computed from the particle locations $\bX_n^i$, and the reconstructed density 
$\rho(t,x)$ is obtained from the Fourier coefficients.

This choice is motivated by the simplicity and spectral accuracy of Fourier methods for smooth periodic functions. However, we note that other functional spaces, such as polynomial bases, wavelets, or kernel-based representations, could offer better adaptability in non-periodic or highly localized settings. We leave the exploration of such alternatives to future work.

\subsection{Extension to PDE system}
To solve PDE systems like Equation \eqref{equ:KS}, we introduce two sets of particles $\bX^i(t)$ and $\bY^i(t)$. Therefore, empirical distribution $\rho_u(t,\intd \bx ) = \mu_x(t,\intd \bx) =  \frac{1}{N_{x,0}}\sum_{i=1}^{N_{x,n}}\delta_{\bX_n^i} $ and $\rho_v(t,\intd \bx ) = \mu_y(t,\intd \bx) =  \frac{1}{N_{y,0}}\sum_{i=1}^{N_{y,n}}\delta_{\bY_n^i}$ are used to approximate $u$ and $v$ separately, by 
\begin{equation}
    u(t,\bx) =  Z_u \rho_u(t,\bx), \quad v(t,\bx) =  Z_v \rho_v(t,\bx),
\end{equation}
where $Z_u = \int_\Omega u(0,\bx) \intd \bx$ and $Z_v = \int_\Omega v(0,\bx) \intd \bx$. Therefore, the dynamics of \eqref{equ:KS} can be written as 
% \begin{equation}
% \begin{aligned}
%     \partial_t u - \nabla \cdot ( \nabla u - \chi(u) \nabla v) &= f_u(u,v),  & x\in \mathbb R^2 ,t\in(0,T),\\
%     v_t -\Delta v & = f_v(u,v),  &x\in \mathbb R^2 , t\in(0,T),
%     \end{aligned}
% \end{equation}
\begin{equation}
    \begin{aligned}
        \partial_t \rho_u =  \mathcal A_u \rho_u + \mathcal B_u \rho_u,\\
        \partial_t \rho_v =  \mathcal A_v \rho_v + \mathcal B_v \rho_v
    \end{aligned}
\end{equation}
where $\mathcal A_u \rho = \nabla \cdot ( \nabla \rho - \frac{Z_v}{Z_u} \chi(Z_u\rho) \nabla \rho_v)$ , $\mathcal A_v \rho = \Delta \rho$, $\mathcal B_u \rho_u  = f_u(Z_u\rho_u,Z_v \rho_v)$ and $B_v \rho_v = f_v(Z_u\rho_u,Z_v \rho_v)$. Similar as before, we can approximate the solution by operator splitting 
\begin{equation}
    \begin{aligned}
        \rho_{u,n}^*  =  \exp(\mathcal A_u\tau) \rho_{u,n}, \quad 
         \rho_{v,n}^*  =  \exp(\mathcal A_v\tau) \rho_{v,n},  \\
         \rho_{u,n+1} =  \exp(\mathcal B_u\tau) \rho_{u,n}^*,  \quad 
         \rho_{v,n+1}  =  \exp(\mathcal B_v\tau) \rho_{v,n}^*,   \\
    \end{aligned}
\end{equation}
where by matching parameters, we have $a_u(\bx;\rho_u,\rho_v) =  - \frac{Z_v}{Z_u\rho_u} \chi(Z_u\rho_u) \nabla \rho_v $ and $a_v(\bx;\rho_u,\rho_v)= 0$. For the diffusion part, we have $\sigma_u(\bx;\rho_u) = \sigma_v(\bx;\rho_v) = 1$. For birth and death part, $\tilde r_u (\bx;\rho_u,\rho_v)=  \frac{f_u(Z_u\rho_u,Z_v \rho_v)}{\rho_u}  $ and $\tilde r_v (\bx;\rho_u,\rho_v )=  \frac{f_v(Z_u\rho_u,Z_v \rho_v)}{\rho_v}  $. The detailed algorithm is shown in  Algorithm \ref{alg:cap}.
\begin{algorithm}
\caption{Stochastic Branching Method for the KS Equation}\label{alg:cap}
\begin{algorithmic}
\Require Initial population size $N_{x,0}, N_{y,0}$, time step size $\tau$, current step $n = 0$, and Fourier mode truncation $K$.
\Require Precomputed normalization constants $Z_u$ and $Z_v$ (either numerically or analytically).
\State Generate $N_{x,0}$ initial particle positions $\bX^i_0$ according to the distribution $\rho_u$ using the Metropolis-Hastings algorithm.
\State Generate $N_{y,0}$ initial particle positions $\bY^i_0$ according to the distribution $\rho_v$ using the Metropolis-Hastings algorithm.
\State Interpolate the particle positions $\bX^i_n$ and $\bY^i_n$ onto Fourier modes up to the truncation $K$ to obtain the Fourier representations $\tilde{\rho}_u$ and $\tilde{\rho}_v$.

\While{$n \tau < T$} \Comment{Assuming $T \mid \tau$}
    \State Update the positions of the particles:
    \For{$i=1$ to $N_{x,n}$}
        \[
        \bX^i_{n,*} = \bX^i_n - \frac{Z_v}{Z_u \tilde{\rho}_u(\bX^i_n)} \chi(Z_u \tilde{\rho}_u(\bX^i_n)) \nabla \tilde{\rho}_v(\bX^i_n) \tau + \sqrt{2\tau} \xi^{i}_{u,n}
        \]
        where $\xi^{i}_{u,n} \sim \mathcal{N}(0, \mathbb{I})$.
    \EndFor
    \For{$i=1$ to $N_{y,n}$}
        \[
        \bY^i_{n,*} = \bY^i_n + \sqrt{2\tau} \xi^{i}_{v,n}
        \]
        where $\xi^{i}_{v,n} \sim \mathcal{N}(0, \mathbb{I})$.
    \EndFor

    \State Interpolate the updated positions $\bX^i_{n,*}$ and $\bY^i_{n,*}$ onto Fourier modes up to truncation $K$ to obtain the Fourier representations $\tilde{\rho}_u^*$ and $\tilde{\rho}_v^*$.

    \For{$i=1$ to $N_{x,n}$}
        \State Compute the branching rate $\tilde{r}_u^i = \frac{f_u(Z_u \rho_u(\bX^i_{n,*}), Z_v \rho_v(\bX^i_{n,*}))}{\rho_u(\bX^i_{n,*})}$.
        \If{$\tilde{r}_u^i > 0$}
            \State Duplicate $\bX^i_{n,*}$ with probability $1 - \exp(-\tilde{r}_u^i \tau)$.
        \ElsIf{$\tilde{r}_u^i < 0$}
            \State Kill $\bX^i_{n,*}$ with probability $1 - \exp(\tilde{r}_u^i \tau)$.
        \EndIf
    \EndFor

    \For{$i=1$ to $N_{y,n}$}
        \State Compute the branching rate $\tilde{r}_v^i = \frac{f_v(Z_u \rho_u(\bY^i_{n,*}), Z_v \rho_v(\bY^i_{n,*}))}{\rho_v(\bY^i_{n,*})}$.
        \If{$\tilde{r}_v^i > 0$}
            \State Duplicate $\bY^i_{n,*}$ with probability $1 - \exp(-\tilde{r}_v^i \tau)$.
        \ElsIf{$\tilde{r}_v^i < 0$}
            \State Kill $\bY^i_{n,*}$ with probability $1 - \exp(\tilde{r}_v^i \tau)$.
        \EndIf
    \EndFor

    \State Update the population sizes: $N_{x,n+1}$ and $N_{y,n+1}$.
    \State Increment the step: $n = n + 1$.
\EndWhile
\end{algorithmic}
\end{algorithm}

% \section{Kolmogorov forward equation and }
% Recall (overdamped) Langevin diffusion is the solution of the following stochastic differential equation,
% where $\bX_t \in \mathbb R^d$ and .  The probability density function $\rho(t)$ of \eqref{equ:overd_langevin} solves the Kolmogorov forward equation on $\mathbb R^d$
% \begin{equation} \label{equ:FPE}
%     \partial_t  \rho =  \nabla \cdot ( \nabla \rho +\rho \bF).
% \end{equation}
% The main idea of the work is to use the branching process to approximate the evolution of the nonlinearity locally. Before combing with the Langevin dynamics, let us consider the pure branching process given by 
% \begin{equation}
%     \partial_t \rho =  V(x) \rho. 
% \end{equation}
% It can be viewed as an infinite system of ordinary differential equations with index $x\in \mathbb R^d$ without any spatial derivative. 
% If $V(t,x) > 0$, then particles located at $x$ undergo duplication, representing local exponential growth. Conversely, if $V(t,x)<0$, particles at 
% $x$ are removed, corresponding to exponential decay. In this way, the pure branching process encodes a pointwise multiplicative effect on the density, without involving any spatial interactions.
% Combining the Langevin dynamics with the branching process yields the following equation for the evolution of the probability density: 
% \begin{equation} \label{equ:FPE_BDL}
% \partial_t \rho = \nabla \cdot ( \nabla \rho + \rho \bF ) + V(x) \rho, 
% \end{equation} 
% where the first term represents diffusion and drift due to the Langevin process, and the second term captures local growth or decay of mass driven by the branching mechanism.

% The solution of \eqref{equ:FPE_BDL} can be approximated by the empirical measure $\mu^{N}_t$ of a collection of interacting particles $\{x_t^i\}_{i=1}^{N_t}$ evolving as followings $\mu^{N}_t = \frac{1}{N} \sum_{i=1}^{N_t} \delta_{x_t^i} $:
% \begin{enumerate}
%     \item[Step 1:] Between birth/death events, each particle $\bX^i_t$ diffuses independently according to \eqref{equ:overd_langevin}.
%     \item[Step 2:] Each particle also has an independent instantaneous birth-death rate $V(\bX_t^i).$
% \end{enumerate}
% Therefore, we can approximate the solution of \eqref{equ:FPE_BDL} by a density estimation.  Unlike the classical mesh-based method, this approach does not require spatial discretization. This is the key reason why this method is promising to overcome the curse of dimensionality in high-dimensional PDEs.


% \begin{proposition}
% Let $\mu^N_t$ be the empirical measure of particles defined above. Assume that $\mu_0^N \rightharpoonup \rho_0$ as $N\to \infty$. Then for all $t\in (0,\infty)$, we have $\mu^N_t \to \rho_t$, where $\rho_t$ solves \eqref{equ:FPE_BDL} with initial condition $\rho_0$
% \end{proposition}
% \begin{proof}
%     First, we notice that the evolution of $\mu^N_t$ is Markovian since the particles $\bX^i_t$ are, and these particles are interchangeable. Following the proof in \cite{rotskoff2019global}, we want to derive the generator and the backward Kolmogorov equation of the measure-valued Markov process. Let's define for some smooth functional $\Psi:\mathcal P  (\mathbb R ^d) \to \mathbb R $, the generator $\mathcal L_N$
%     \[
%     (\mathcal L_N \Psi ) (\mu^N) := \lim_{t\to 0^+} \frac{\mathbb E_0 \Psi ( \mu_t^N) - \Psi(\mu^N)}{t},
%     \]
%     where $\mathbb E_0 $ denotes the expectation of $\Psi(\mu^N_t)$ conditioned on that $\mu_0^N=\mu^N$. To compute the limit, notice that if particle $\bX_t^i$ is killed at time $t$, the changes this indeuced on $\mu^N_t$ is 
%     \[
%     \mu^N_t =  \mu_{t_-}^N - \frac{1}{N} \delta_{\bX^i}.
%     \]
%     Similarly if particle $\bX_t^i$ is duplicated at time $t$, the changes this indeuced on $\mu^N_t$ is 
%     \[
%     \mu^N_t =  \mu_{t_-}^N + \frac{1}{N} \delta_{\bX^i}.
%     \]
%     If we account for the rate at which these events occur, as well as the effect of the diffusion process, we can explicitly compute the generator 
%     \[
%     \begin{aligned}
%     (\mathcal L _N \Psi)(\mu^N) =& \frac{1}{N} \sum_i ^N \int \left( 
%     \bF(\bX^i) \cdot \nabla D_{\mu^N} \Psi(\bX^i)  
%     + \Delta D_{\mu^N} \Psi(\bX^i)
%     \right) \delta_{\bX^i}(\intd x) \\
%     &+ \frac{1}{N} \sum_{i=1}^N \int V(\bX^i) \delta_{\bX^i}(\intd x)\left(\Psi(\mu^N + \frac{1}{N} \text{sign} V(\bX^i)\delta_{\bX^i} ) - \Psi(\mu^N)\right),
%     \end{aligned}
%     \]
%     where the functional derivative $D_\mu \Psi$ is the function from $\mathbb R^d \to \mathbb R$ defined by that for any signed measure $\nu$ with $\int \nu(\intd x)= 0  $, 
%     \[
%     \lim_{\epsilon \to 0 } \frac{\Psi(\mu+\epsilon\nu)-\Psi(\mu)}{\epsilon} = \int D_\mu\Psi(x)\nu(\intd x).
%     \]
%     By the definition of the empirical measure $\mu^N$, the generator $\mathcal L_N$ can be rewritten as 
%     \[
%     \begin{aligned}
%     (\mathcal L _N \Psi)(\mu^N) =& \int \left( 
%     \bF(x) \cdot \nabla D_{\mu^N} \Psi(x)  
%     + \Delta D_{\mu^N} \Psi(x)
%     \right) \mu^N(\intd x)  \\
%     +&  \int V(x) \mu^N(\intd x) \left(\Psi(\mu^N + \frac{1}{N} \text{sign} V(x)\delta_{x} ) - \Psi(\mu^N)\right) ,
%     \end{aligned}
%     \]
% \end{proof}
% \section{Stochastic Branching method}
% In this work, we aim to extend the above framework to fully nonlinear PDEs. A natural strategy is to adopt an iterative scheme that treats the nonlinear equation as a linear one at each iteration, conditioned on the current estimate of the solution. Specifically, at each time $t$, the equation \eqref{equ:KS} can be approximated by the following linearized system: 
% \begin{equation}
% \begin{aligned}
%     \partial_t \rho_u &=  \nabla \cdot \left ( \bF_u(t,x) \rho_u +\nabla \rho_u \right) + V_u(t,x) \rho_u,\\
%     \partial_t \rho_v &=  \nabla \cdot \left ( \bF_v(t,x) \rho_v +\nabla \rho_v \right) + V_v(t,x) \rho_v,
% \end{aligned}
% \end{equation}
% where the drift and reaction terms are given by $\bF_u =  -\frac{\chi(u)}{u}\nabla v, \bF_v = 0, V_u = \frac{f_u(u,v)}{u}$ and $V_v= \frac{f_v(u,v)}{v}$. At each iteration, the fields $u$ and $v$ are reconstructed from the particle densities as  $u = C_u \rho_u$ and $v= C_v\rho_v$, where $C_u = \int u(0,\bx)\intd \bx$ and $C_v = \int v(0,\bx)\intd \bx$ are normalization constants.

\section{Numerical Examples}
In this section, we present numerical examples in two-dimensional space to verify the accuracy of our method.


\textbf{Case 1:} We begin by solving the following problem in a periodic domain $\Omega = [0, 2\pi] \times [0, 2\pi]$:
\begin{equation} \label{equ:case1}
\begin{aligned}
&u_t - \nabla \cdot \left(\nabla u - u \nabla v \right) = 0  , &\quad \bx \in \Omega , t > 0, \\
&v_t - \Delta v = u - v , & \quad \bx \in \Omega, t > 0.
\end{aligned}
\end{equation}
The system is subject to periodic boundary conditions, with initial conditions given by:
$$
u_0 (\bx) =  \sin^2(\bx_1) \cos^2(\bx_2), \quad v_0(\bx) =  \cos(\bx_1) + \cos(\bx_2) + 2.
$$
The temporal evolution of the total mass of $v$ is characterized by
$$
M(t) = \int v(t,\bx) \, \intd \bx.
$$
It can be shown that the rate of change of the mass is given by
$$
\dot{M}(t) = -M(t) + \int u(x,0) \, dx = -M(t) + \pi^2.
$$
We present the numerical reuslts for $M(t)$, comparing our method with the exact solution in Figure~\ref{fig:total_mass}.

\begin{figure}[h]
\centering
\includegraphics[width=0.8\linewidth]{figure/total_mass.pdf}
\caption{Time evolution of the total mass of $v$ for different numbers of particles. The left panel shows the total mass over time, while the right panel displays the corresponding error in the total mass.}
\label{fig:total_mass}
\end{figure}

Additionally, we compare our particle-based method with a finite difference method with a grid of $100 \times 100$ points. The result for $u(t)$ and $v(t)$ at $t = 0.0, 0.1, 0.2, 0.3, 0.4$ are shown in Figures~\ref{fig:fd_comparison_u} and~\ref{fig:fd_comparison_v}, respectively. Figure~\ref{fig:convergence_error} illustrates the reduction in relative $L^2$ error as the number of samples increases, showing a clear $1/\sqrt{N}$ convergence rate.

\begin{figure}[h]
\centering
\includegraphics[width=\linewidth]{figure/case1_u.pdf}
\caption{Comparison of $u$ between the finite difference method with a $100 \times 100$ grid and our particle-based method at $t = 0.0, 0.1, 0.2, 0.3, 0.4$.}
\label{fig:fd_comparison_u}
\end{figure}

\begin{figure}[h]
\centering
\includegraphics[width=\linewidth]{figure/case1_v.pdf}
\caption{Comparison of $v$ between the finite difference method with a $100 \times 100$ grid and our particle-based method at $t = 0.0, 0.1, 0.2, 0.3, 0.4$.}
\label{fig:fd_comparison_v}
\end{figure}

\begin{figure}[h]
\centering
\includegraphics[width=0.8\linewidth]{figure/case1_err.pdf}
\caption{Relative $L^2$ error of $u$ and $v$ with respect to the number of samples $N_0$, compared with $N_0 = 320,000$.}
\label{fig:convergence_error}
\end{figure}

\textbf{Case 2:} For the second case, we solve the system \eqref{equ:case1} with the initial conditions:

$$
u_0(\bx) = 840 \exp\left(-84 (\bx_1^2 + \bx_2^2)\right), \quad v_0(\bx) = 420 \exp\left(-42(\bx_1^2 + \bx_2^2)\right).
$$

The result for $u$ is shown in Figure~\ref{fig:case2_u}, where we observe that the solution undergoes a blow-up, a phenomenon that is generally difficult to capture using traditional function spaces. Furthermore, the $L^2$ norm of the numerical solution for different values of $N_0$ is shown in Figure~\ref{fig:case2_l2}, demonstrating that varying $N_0$ results in similar growth patterns, in contrast to the discontinuous Galerkin method \cite{li2017local}, where the solution is highly sensitive to the mesh size.

\begin{figure}
\centering
\includegraphics[width=0.8\linewidth]{figure/case2_u.pdf}
\caption{Numerical approximations of $u$ at $t = 0, 5 \times 10^{-5}, 10^{-4}, 1.5 \times 10^{-4}$.}
\label{fig:case2_u}
\end{figure}

\begin{figure}
\centering
\includegraphics[width=0.5\linewidth]{figure/case2_L2.pdf}
\caption{Time evolution of the $L^2$-norm of the numerical approximation for different values of $N_0$.}
\label{fig:case2_l2}
\end{figure}

\textbf{Case 3:} Lastly, we examine a fully nonlinear case with logistic growth for the cells, where the interaction between the particles is not known a priori. The system is given by:
\begin{equation}
\label{equ:nonliear}
\begin{aligned}
&u_t - \nabla \cdot \left(\nabla u - \frac{4u}{1 + u^2} \nabla v \right) = u(1 - u), & \quad \bx \in \Omega, t > 0, \\
&v_t - \Delta v = u - v, & \quad \bx \in \Omega, t > 0.
\end{aligned}
\end{equation}
The theoretical results of this model can be found in \cite{osaki2002exponential, wang2007classical, wrzosek2004global}. The initial conditions are specified as:
$$
    \begin{aligned}
        u_0(\bx) =&  \exp\left( (\bx_1 - 0.5\pi)^2 + (\bx_2 - 0.5\pi)^2 \right) + \exp\left( (\bx_1 - \pi)^2 + (\bx_2 - \pi)^2 \right) \\
        &+ \exp\left( (\bx_1 - 1.5\pi)^2 + (\bx_2 - 1.5\pi)^2 \right), \\
        v_0(\bx) =&  \exp\left( (\bx_1 - 0.4\pi)^2 + (\bx_2 - 0.4\pi)^2 \right) + \exp\left( (\bx_1 - 0.8\pi)^2 + (\bx_2 - 0.8\pi)^2 \right) \\
        &+ \exp\left( (\bx_1 - 1.2\pi)^2 + (\bx_2 - 1.2\pi)^2 \right) + \exp\left( (\bx_1 - 1.6\pi)^2 + (\bx_2 - 1.6\pi)^2 \right).
    \end{aligned}
$$
Since both the birth and death of cells, as well as the chemical concentrations $f_u$ and $f_v$, are non-zero, the particle method proposed in \cite{wang2025novel}, which only accounts for advection and diffusion, has limited capability in addressing these effects.  In contrast, our method is sufficiently general to handle these additional complexities.
Figures~\ref{fig:case3_u} and~\ref{fig:case3_v} display the time evolution of the solutions $u$ and $v$, respectively, obtained using both the particle-based and finite difference methods.

\begin{figure}
\centering
\includegraphics[width=0.8\linewidth]{figure/case3_u.pdf}
\caption{Time evolution of the solution $u$ computed using the particle-based and finite difference methods.}
\label{fig:case3_u}
\end{figure}

\begin{figure}
\centering
\includegraphics[width=0.8\linewidth]{figure/case3_v.pdf}
\caption{Time evolution of the solution $v$ computed using the particle-based and finite difference methods.}
\label{fig:case3_v}
\end{figure}

\section{Discussion}
In this work, we introduce a novel framework for solving a class of reaction-diffusion-advection equations. Unlike existing particle-based or Lagrangian methods \cite{qiu2010conservative, qiu2011positivity, wang2025novel, carrillo_interacting_2024}, our approach is more general and extends beyond conservation laws, utilizing a purely particle-based scheme. The only place where we incorporate a Fourier basis is in the nonlinear interpolation, which could potentially be replaced by blob methods~\cite{campos2018convergence,craig2016blob,goodman1990convergence,hauray2009wasserstein} for high-dimensional problems in the future.

High-dimensional PDEs have long posed a significant challenge. Recent advances in neural network-based methods, such as Deep Galerkin Methods (DGM) \cite{sirignano2018dgm}, Deep Ritz \cite{yu2018deep}, and MIM \cite{lyu2022mim}, have faced limitations in accuracy due to the inherent difficulties of nonlinear optimization. Similarly, methods like Deep-BSDE \cite{han2018solving} and the Deep Picard approach \cite{han2024deep} also rely on particle or SDE-based techniques. However, nonlinear optimization remains a critical bottleneck in these methods. Looking ahead, we are particularly interested in exploring Fourier interpolation using low-rank techniques~\cite{tang2025solving} to tackle high-dimensional problems more efficiently.




\bibliographystyle{unsrt}
\bibliography{main}  %%% Remove comment to use the external .bib file
\clearpage

\end{document}
